\section{Special Functions}
\label{sec:specialfns}

The elementary functions---polynomials, the exponential, the trigonometric and
their inverses---run out long before mathematics does. The moment you write down
an integral like \(\int_0^\infty t^{z-1}e^{-t}\,dt\), or a power series whose
coefficients grow by a fixed ratio, or the solution of a second-order linear
differential equation with a variable coefficient, you have left the elementary
world and entered the realm of the \emph{special functions}: the gamma and zeta
functions, the error function, the Bessel and Airy functions, the Legendre
functions and the classical orthogonal polynomials, the elliptic integrals, the
hypergeometric family up to Meijer's \(G\), and the Dirichlet
\(L\)-functions of number theory.\index{special functions} They are ``special'' only in
the sense that they were named and tabulated before computers existed; today they
are as ordinary as \(\sin\), and Mathilda treats them so.

What makes a special function useful in a computer algebra system is that it wears
three faces at once, and Mathilda shows you all three. There are the
\emph{exact symbolic values}---\(\Gamma(1/2)=\sqrt\pi\), \(\zeta(2)=\pi^2/6\)---which a
numeric library can only approximate. There are the \emph{symbolic identities}: the
recurrences, reflection and derivative rules, series expansions, and orthogonality
relations that let these functions participate in algebra rather than merely
returning numbers. And there is \emph{rigorous arbitrary-precision numerics}, for
when a number really is what you want, computed to any precision you ask for and,
where the kernel is backed by FLINT's ball arithmetic, with a guaranteed error
bound. This section takes the major families in turn, and the same three faces
recur in each: what closes into a clean symbol, what identity the system knows, and
how it puts a number on the value.

\subsection{The gamma family}
\label{ssec:sf-gamma}

Everything begins with the \B{Gamma} function, \(\Gamma(z)=\int_0^\infty
t^{z-1}e^{-t}\,dt\), the continuous interpolation of the
factorial.\index{gamma function} Its defining property, \(\Gamma(z+1)=z\,\Gamma(z)\)
with \(\Gamma(1)=1\), forces \(\Gamma(n)=(n-1)!\) at the integers, and from the
single exact value \(\Gamma(1/2)=\sqrt\pi\) the recurrence generates the whole
ladder of half-integers.

\begin{codepairs}
\pair{special-functions/gamma/1}{At the positive integers \B{Gamma} is the shifted
factorial: \(\Gamma(n)=(n-1)!\), so this is \(0!\) through \(5!\).}
\pair{special-functions/gamma/2}{The one transcendental seed value,
\(\Gamma(1/2)=\sqrt\pi\), returned exactly.}
\pair{special-functions/gamma/3}{The recurrence climbs from it:
\(\Gamma(7/2)=\tfrac{15}{8}\sqrt\pi\), a rational multiple of \(\sqrt\pi\).}
\pair{special-functions/gamma/4}{And descends to the negative half-integers---here
\(\Gamma(-1/2)=-2\sqrt\pi\), finite because the poles fall only on the
\emph{integers}.}
\pair{special-functions/gamma/5}{At the non-positive integers \(\Gamma\) has simple
poles, reported as \mcode{ComplexInfinity}.}
\pair{special-functions/gamma/6}{The two-argument form is the upper incomplete
gamma \(\Gamma(a,z)\); at a positive integer \(a\) it closes into elementary
form.}
\pair{special-functions/gamma/7}{With a numeric second argument that closed form is
an exact number: \(\Gamma(5,2)=168/e^{2}\).}
\pair{special-functions/gamma/8}{Where no symbol exists, \B{N} delivers digits---here
\(\Gamma(i)\) to thirty places, real and imaginary parts together.}
\pair{special-functions/gamma/9}{And \(\Gamma\) knows its own derivative, which
closes through the digamma function \(\psi=\Gamma'/\Gamma\).}
\end{codepairs}

The half-integer ladder (pairs~2--4) is the gamma function's signature, and it is
worth seeing why it terminates in \(\sqrt\pi\) rather than another special value.
\calloutnote{Theory}{The reflection formula
\(\Gamma(z)\,\Gamma(1-z)=\pi/\sin(\pi z)\)\index{gamma function!reflection formula}
evaluated at \(z=1/2\) gives \(\Gamma(1/2)^2=\pi\), hence \(\Gamma(1/2)=\sqrt\pi\).
Combined with the recurrence \(\Gamma(z+1)=z\,\Gamma(z)\), it pins every
half-integer to a rational multiple of \(\sqrt\pi\). The same reflection formula
explains the pole structure: since \(1/\Gamma\) is entire with zeros exactly at
\(0,-1,-2,\dots\), \(\Gamma\) itself has simple poles there and nowhere else, which
is why pair~4 is finite but pair~5 is not.} The numeric face (pairs~7--8) is the
one a differential-equation solver or a statistics routine actually leans
on.\calloutnote{Under the hood}{\B{Gamma} lives in
\texttt{src/special\_functions/gamma.c}. A machine-precision real routes to the C
library's \texttt{tgamma}; a machine-precision \emph{complex} argument to Mathilda's own
double-precision kernel (\texttt{sf\_machine\_gamma.c}), which reads \(\log\Gamma\) off
precomputed Taylor tiles near the real axis and Stirling's series beyond, and is the same
code the packed-array and \B{Compile} paths call, so all three agree to the bit; an
arbitrary-precision real goes to MPFR's correctly rounded \texttt{mpfr\_gamma}; and an
arbitrary-precision complex argument to Spouge's
approximation, whose coefficients are recomputed at runtime to the precision you
request, with the reflection formula handling \(\operatorname{Re}(z)<\tfrac12\). The
incomplete form uses \texttt{mpfr\_gamma\_inc} on the real axis and a lower-incomplete
series or a Lentz continued fraction in the complex plane.}

\usagebox{Gamma}

The \B{Beta} function \(B(a,b)=\Gamma(a)\Gamma(b)/\Gamma(a+b)\) is the gamma
function's closest relative---the normalizing constant of the beta distribution and
the value of the integral \(\int_0^1 t^{a-1}(1-t)^{b-1}\,dt\).\index{beta function}
Because it is built entirely from gammas, its exact reductions follow from theirs.

\begin{codepairs}
\pair{special-functions/beta/1}{A positive integer argument collapses the gamma
ratio to a rational: \(B(2,3)=1/12\).}
\pair{special-functions/beta/2}{Even with a rational partner it stays exact,
\(B(3,1/3)=27/14\).}
\pair{special-functions/beta/3}{Half-integers fold to rational multiples of \(\pi\)
instead---\(B(1/2,1/2)=\pi\), the arcsine normalization.}
\pair{special-functions/beta/4}{And \(B(5/2,7/2)=\tfrac{3}{256}\pi\).}
\pair{special-functions/beta/5}{When no such collapse is available, the residual
gamma form survives rather than being forced to a number.}
\pair{special-functions/beta/6}{A fully numeric call evaluates like any other.}
\end{codepairs}

Pair~5 is a small lesson in the system's temperament: \(B(1/3,1/3)\) is left as
\(\Gamma(1/3)^2/\Gamma(2/3)\) because no further exact simplification is available
without the reflection formula, which Mathilda does not apply automatically to a
product of gammas at complementary rationals.\calloutnote{Pitfall}{Do not expect
\(\Gamma(1/3)\,\Gamma(2/3)\) to collapse to \(2\pi/\sqrt3\) on its own. The gamma
reflection and multiplication identities are known to the theory but are \emph{not}
auto-applied---applying them blindly would fight as often as it helps---so products
of gammas at rational points are carried symbolically until you rewrite them by
hand. This is the honest cost of not over-simplifying.}

Two functions read the shape of the gamma function near its poles and along the
real axis. \B{LogGamma} is the log-gamma function \(\log\Gamma(z)\): the analytic
continuation of \(\log(\Gamma(z))\), not the same thing as \(\log(\Gamma(z))\)
term by term, and the right primitive whenever \(\Gamma\) itself would overflow.
\B{PolyGamma} is the family of logarithmic derivatives: \(\psi(z)=\psi^{(0)}(z)=
\Gamma'(z)/\Gamma(z)\) is the digamma function,\index{digamma function} and
\(\psi^{(n)}\) its \(n\)-th derivative.

\begin{codepairs}
\pair{special-functions/polygamma/1}{\B{LogGamma} keeps integers exact as the log of
a factorial: \(\log\Gamma(10)=\log(362880)=\log(9!)\).}
\pair{special-functions/polygamma/2}{At \(z=171\) the gamma function has already
overflowed a machine double---}
\pair{special-functions/polygamma/3}{---but \B{LogGamma} sails on, returning
\(\approx706.6\). This is why factorial ratios and log-likelihoods are computed
through it, never through \(\Gamma\) directly.}
\pair{special-functions/polygamma/4}{The digamma at a positive integer is a rational
minus Euler's constant: \(\psi(5)=\tfrac{25}{12}-\gamma\).}
\pair{special-functions/polygamma/5}{Gauss's digamma theorem closes rational
arguments into elementary form: \(\psi(1/2)=-\gamma-\log 4\).}
\pair{special-functions/polygamma/6}{The trigamma at \(1\) is \(\zeta(2)=\pi^2/6\).}
\pair{special-functions/polygamma/7}{\B{LogGamma} differentiates to digamma.}
\pair{special-functions/polygamma/8}{Each further derivative raises the
polygamma order by one.}
\end{codepairs}

The contrast between pairs~2 and~3 is the entire reason \B{LogGamma} exists as a
first-class head rather than a shorthand for \texttt{Log[Gamma[z]]}: it stays finite
across a range where the gamma function runs from \(10^{-300}\) to \(10^{300}\), and
its imaginary part tracks the \emph{continued} argument past \(\pi\), where a naive
\(\log(\Gamma(z))\) would wrap.\calloutnote{Under the hood}{\B{PolyGamma}
(\texttt{src/special\_functions/polygamma.c}) handles order \(n=0\) at rational
arguments through the Gauss digamma theorem---a finite trigonometric sum---and closes
odd orders through \(\zeta(n{+}1)\). Numerically it uses MPFR's digamma for
\(n=0\) and a recurrence-shift plus a Bernoulli asymptotic expansion for higher
orders; complex arbitrary-precision values are handed to FLINT's \texttt{acb\_polygamma}.
The negative order \(\psi^{(-1)}(z)\) is \B{LogGamma} itself.}

The last member of the family is the \B{Pochhammer} symbol \((a)_n=
a(a+1)\cdots(a+n-1)=\Gamma(a+n)/\Gamma(a)\), the \emph{rising factorial} that
supplies the coefficients of every hypergeometric series we meet at the end of this
section.\index{Pochhammer symbol}

\begin{codepairs}
\pair{special-functions/pochhammer/1}{A symbolic base gives an explicit polynomial
product of \(n\) linear factors.}
\pair{special-functions/pochhammer/2}{A numeric base gives an exact value,
\((10)_6=3603600\).}
\pair{special-functions/pochhammer/3}{\((1)_{25}\) is just \(25!\), carried exactly
as a GMP big integer.}
\pair{special-functions/pochhammer/4}{Half-integer arguments reduce through the gamma
ratio: \((1/2)_3=15/8\).}
\pair{special-functions/pochhammer/5}{And \((3/2)_{1/2}=2/\sqrt\pi\), a gamma ratio
that happens to close.}
\end{codepairs}

\subsection{The zeta family}
\label{ssec:sf-zeta}

The Riemann zeta function \(\zeta(s)=\sum_{n\ge1}n^{-s}\) is the deepest object in
this section, and \B{Zeta} carries its exact arithmetic.\index{Riemann zeta
function} At the even positive integers it closes into powers of \(\pi\); its values
at the non-positive integers, defined by analytic continuation, are rational; and
its values at the odd positive integers are believed to be new transcendentals with
no closed form at all.

\begin{codepairs}
\pair{special-functions/zeta/1}{The even values are rational multiples of powers of
\(\pi\): \(\zeta(2)=\pi^2/6\).}
\pair{special-functions/zeta/2}{The next even value is \(\zeta(4)=\pi^4/90\).}
\pair{special-functions/zeta/3}{This is exactly the Basel sum, which \B{Sum}
recognizes independently and returns identically.}
\pair{special-functions/zeta/4}{Analytic continuation gives \(\zeta(0)=-1/2\).}
\pair{special-functions/zeta/5}{It also gives the famous \(\zeta(-1)=-1/12\).}
\pair{special-functions/zeta/6}{The negative even integers are the \emph{trivial
zeros}: \(\zeta(-2)=0\).}
\pair{special-functions/zeta/7}{The odd values have no known closed form, so
\(\zeta(3)\) is returned as itself---a genuine constant, not a failure.}
\pair{special-functions/zeta/8}{But it is a \emph{number}: \B{N} gives Apéry's
constant to thirty digits.}
\pair{special-functions/zeta/9}{On the critical line, at the height \(t=14.1347\ldots\)
of the first nontrivial zero, \(\zeta(1/2+it)\) collapses to the order of
\(10^{-7}\)---the Riemann hypothesis made visible.}
\pair{special-functions/zeta/10}{The Hurwitz zeta \(\zeta(s,a)=\sum_{n\ge0}(n+a)^{-s}\)
shifts the denominators, and still closes: \(\zeta(2,3)=\pi^2/6-5/4\).}
\pair{special-functions/zeta/11}{The Stieltjes constants are inert symbols; the
zeroth is Euler's constant \(\gamma\).}
\end{codepairs}

Pairs~4--6 are not sleight of hand but the content of the \emph{functional
equation}, and pair~9 is a numerical glimpse of the most famous open problem in
mathematics.\calloutnote{Theory}{Riemann's functional
equation\index{Riemann zeta function!functional equation}
\(\zeta(s)=2^s\pi^{s-1}\sin(\pi s/2)\,\Gamma(1-s)\,\zeta(1-s)\) relates the value at
\(s\) to the value at \(1-s\), and with it the whole left half-plane is determined by
the right. The factor \(\sin(\pi s/2)\) vanishes at the negative even integers,
producing the trivial zeros of pair~6; the values \(\zeta(0)=-1/2\) and
\(\zeta(-1)=-1/12\) fall out of the same equation. The \emph{nontrivial} zeros all
lie in the strip \(0<\operatorname{Re}(s)<1\), and the \emph{Riemann hypothesis}
\index{Riemann hypothesis} asserts they lie exactly on the line
\(\operatorname{Re}(s)=1/2\); pair~9 evaluates \(\zeta\) at the first of them, where
it is zero to the precision requested.} The even values, by contrast, are elementary
once you know the Bernoulli numbers.\calloutnote{Theory}{Euler's formula
\(\zeta(2n)=(-1)^{n+1}\dfrac{B_{2n}(2\pi)^{2n}}{2\,(2n)!}\) expresses every even zeta
value through a \B{BernoulliB} number\index{Bernoulli numbers}; for \(n=1\),
\(B_2=1/6\) gives \(\zeta(2)=\pi^2/6\). No analogous formula is known for the odd
values, which is why \(\zeta(3)\) (pair~7) has no closed form while \(\zeta(2)\) and
\(\zeta(4)\) do.} Underneath the exact layer is a rigorous numeric
one.\calloutnote{Under the hood}{\B{Zeta}, \B{HurwitzZeta} and \B{StieltjesGamma}
(\texttt{src/special\_functions/zeta.c}, and \texttt{src/flint\_num\_bridge.c} for
the numerics) delegate arbitrary-precision \emph{complex} evaluation to FLINT's Arb
library---\texttt{acb\_dirichlet\_zeta}, \texttt{acb\_dirichlet\_hurwitz},
\texttt{acb\_dirichlet\_stieltjes}. Arb works in \emph{ball arithmetic}\index{ball
arithmetic}: every value carries a midpoint and a rigorous radius, so pair~9's
near-zero comes with a proof that the true value lies within a known ball of the
printed digits---not a hope that rounding behaved.}

\usagebox{Zeta}

The Stieltjes constants \(\gamma_n\) (pair~11) are the coefficients of the Laurent
expansion of \(\zeta\) about its pole at \(s=1\); the zeroth is Euler's constant
\(\gamma\), surfaced here as \B{EulerGamma}, and the rest are carried as inert
\B{StieltjesGamma} symbols.

The Hurwitz zeta of pair~10 is defined for \(a>0\), but the sum makes sense for any
offset that avoids the one term \(k+a=0\). For \(a\le 0\) Mathilda follows the
reference system's convention, summing \(\bigl((k+a)^2\bigr)^{-s/2}\) so that a
negative base never forces a branch choice and dropping the singular term; it then
shifts the offset into \((0,1]\) by peeling off the first few terms exactly,
\(\zeta(s,a)=\sum_{k<m,\;k+a\ne0}|k+a|^{-s}+\zeta(s,a+m)\), and closes the remainder
when it can.\index{Riemann zeta function!Hurwitz, negative offset}

\begin{codepairs}
\pair{special-functions/zeta-hurwitz/1}{One peeled term, \((1/2)^{-2}=4\), plus
\(\zeta(2,1/2)=(2^2-1)\zeta(2)=\pi^2/2\).}
\pair{special-functions/zeta-hurwitz/2}{Two peeled terms, \(4/9\) and \(4\), on top of
the same remainder.}
\pair{special-functions/zeta-hurwitz/3}{At \(s=-1\) the peeled terms are \(2+1\), and
\(\zeta(-1,1)=-1/12\) finishes the job.}
\pair{special-functions/zeta-hurwitz/4}{At an integer offset the \(k+a=0\) term is
dropped: \(1^{-2}+\zeta(2)\).}
\pair{special-functions/zeta-hurwitz/5}{\mcode{IncludeSingularTerm -> True} keeps it,
and \(0^{-2}\) is a pole.}
\end{codepairs}

The machine path agrees: a real \(a\le0\) runs the same peel-and-shift in double
precision, with the peeled terms summed compensatedly, and it is the same kernel that
a packed array of offsets or a compiled function calls.

\subsection{Dirichlet characters and \texorpdfstring{\(L\)}{L}-functions}
\label{ssec:sf-dirichlet}

Euler's product \(\zeta(s)=\prod_p(1-p^{-s})^{-1}\) encodes the primes; Dirichlet's
insight was that twisting it by a periodic sign pattern separates the primes by
residue class. A \emph{Dirichlet character} modulo \(k\) is a function
\(\chi:\mathbb Z\to\mathbb C\) that is periodic mod \(k\), completely multiplicative,
and zero exactly on the integers sharing a factor with \(k\); its nonzero values are
roots of unity, and there are \(\varphi(k)\) of
them.\index{Dirichlet character} Its \(L\)-function is the twisted zeta
series\index{Dirichlet L-function}
\[
  L(s,\chi)=\sum_{n\ge1}\frac{\chi(n)}{n^{s}}=\prod_p\bigl(1-\chi(p)\,p^{-s}\bigr)^{-1}.
\]
\B{DirichletCharacter}\texttt{[k, j, n]} gives \(\chi(n)\) for the \(j\)-th character
mod \(k\), and \B{DirichletL}\texttt{[k, j, s]} its \(L\)-function; \(j=1\) is
always the principal character. The number theory that uses them is
\S\ref{sec:numbertheory}; here they are special functions, with the same three faces
as \(\zeta\).

\begin{codepairs}
\pair{special-functions/dirichlet/1}{The one non-principal character mod 4:
\(+1\) on \(n\equiv1\), \(-1\) on \(n\equiv3\), \(0\) on the evens.}
\pair{special-functions/dirichlet/2}{Mod 5 there are four characters; this one has
order 4, so its values are the fourth roots of unity.}
\pair{special-functions/dirichlet/3}{The real character mod 5 is the Legendre symbol
\(\bigl(\frac{n}{5}\bigr)\), checked against \B{JacobiSymbol}.}
\pair{special-functions/dirichlet/4}{Leibniz's series \(1-\tfrac13+\tfrac15-\cdots\)
is \(L(1,\chi_4)=\pi/4\).}
\pair{special-functions/dirichlet/5}{Its mod-3 analogue, \(\pi/(3\sqrt3)\).}
\pair{special-functions/dirichlet/6}{\(L\) is rational at every non-positive
integer; for instance \(L(-2,\chi_4)=-1/2\).}
\pair{special-functions/dirichlet/7}{The principal character mod 1 is the constant
\(1\), so \(L\) is \(\zeta\) itself.}
\pair{special-functions/dirichlet/8}{Mod 4 the principal character removes the prime
2: \(L(s,\chi_0)=(1-2^{-s})\zeta(s)\), here \((1-2)(-\tfrac1{12})\).}
\pair{special-functions/dirichlet/9}{\(L(2,\chi_4)\) has no closed form in
\mcode{DirichletL}, but its digits are---}
\pair{special-functions/dirichlet/10}{---Catalan's constant \(G\), \B{Catalan}.}
\pair{special-functions/dirichlet/11}{The real character mod 5 at \(s=1\) comes out
as a difference of logarithms of sines.}
\pair{special-functions/dirichlet/12}{To thirty digits that difference is
\(2\log\phi/\sqrt5\), with \(\phi\) the golden ratio.}
\pair{special-functions/dirichlet/13}{A machine value on the critical line,}
\pair{special-functions/dirichlet/14}{the same point to forty digits,}
\pair{special-functions/dirichlet/15}{and a complex character at a complex \(s\).}
\pair{special-functions/dirichlet/16}{An \B{Interval} argument returns a rigorous
enclosure: \(L(s,\chi_4)\) runs from \(G\) at \(s=2\) to \(\pi^3/32\) at \(s=3\).}
\end{codepairs}

The exact values come from two classical theorems.\calloutnote{Theory}{For a
non-principal \(\chi\) the series converges at \(s=1\), and \(L(1,\chi)\) is a finite
sum of logarithms (an even character) or of \(\pi\) times algebraic numbers (an odd
one)---the input to Dirichlet's class number formula and the reason \(L(1,\chi)\ne0\),
which is the heart of his theorem on primes in arithmetic progressions. At the
non-positive integers, \(L(1-n,\chi)=-B_{n,\chi}/n\), where \(B_{n,\chi}\) is a
\emph{generalized Bernoulli number}\index{Bernoulli numbers!generalized}
\(B_{n,\chi}=k^{n-1}\sum_{a=1}^{k}\chi(a)B_n(a/k)\); this is the twisted form of
\(\zeta(1-n)=-B_n/n\) and why pair~6 is rational. Like \(\zeta\), every primitive
\(L(s,\chi)\) satisfies a functional equation relating \(s\) to \(1-s\), and the
generalized Riemann hypothesis puts its nontrivial zeros on
\(\operatorname{Re}s=\tfrac12\).} Two things need care in a CAS, and both are about
\emph{which} character \(j\) names.\calloutnote{Under the hood}{There is no canonical
numbering of the characters mod \(k\). Mathilda uses the reference system's \((k,j)\)
labelling: the primes of \(k\) in ascending order, each odd prime power generated by
its smallest primitive root, \(2^a\) (\(a\ge3\)) split into \(\langle-1\rangle\) then
\(\langle5\rangle\), and \(j-1\) read as a mixed-radix number with the first
component most significant. The convention is not published; it was reconstructed and
verified on 148{,}138 character values. Evaluating \(\chi(n)\) needs a discrete
logarithm in each component, done by Pohlig--Hellman. Numerically, arbitrary
precision goes to FLINT's \texttt{acb\_dirichlet\_l}, with the character set on
FLINT's \emph{own} generators rather than translated through its (different)
labelling; machine precision has a native kernel (\texttt{sf\_machine\_dirichlet.c})
that runs Euler--Maclaurin on \(L\) with the integral terms written through
\texttt{expm1}/\texttt{log1p}, so nothing cancels as \(s\to1\), and the functional
equation for \(\operatorname{Re}s<\tfrac12\). The engine is
\texttt{src/special\_functions/dirichlet.c}.}

\usagebox{DirichletL}

Pair~11 shows the system's habit of returning the closed form it derived rather than
the one you know: the class-number formula naturally produces logarithms of
\(\sin(\pi a/5)\), and recognizing their ratio as the golden ratio is a separate
simplification. The 30-digit comparison in pair~12 is the quick way to confirm such an
identity.\calloutnote{Pitfall}{Do not take machine \(L\)-values from another system
on trust in the left half-plane, where the functional equation multiplies by large
gamma factors. The special-functions benchmark (\S\ref{ssec:sf-machine}) found the
reference system's own machine value at \(L(-3.3,\chi)\), \(\chi\) the real character
mod 5, to be \(2.492235964\) where the true value is \(2.4922356398\ldots\)---wrong from
the seventh digit. When an \(L\)-value matters, ask for it at an exact \(s\) with
\mcode{N[..., 20]}, which in Mathilda is evaluated in ball arithmetic.}

\subsection{Error functions and the integral functions}
\label{ssec:sf-erf}

The error function \(\operatorname{erf}(z)=\frac{2}{\sqrt\pi}\int_0^z e^{-t^2}\,dt\)
is the antiderivative of the Gaussian, and no elementary function equals
it---which is precisely why it needs a name.\index{error function} \B{Erf}, its
complement \B{Erfc}, the imaginary variant \B{Erfi}, and the inverse \B{InverseErf}
form a small closed world.

\begin{codepairs}
\pair{special-functions/erf/1}{At the origin \(\operatorname{erf}(0)=0\).}
\pair{special-functions/erf/2}{At infinity \(\operatorname{erf}(\infty)=1\): the total
mass of the normalized Gaussian.}
\pair{special-functions/erf/3}{It is an odd function, and the odd symmetry is applied
symbolically.}
\pair{special-functions/erf/4}{Arbitrary-precision numerics track the digits you ask
for.}
\pair{special-functions/erf/5}{The complement \(\operatorname{erfc}=1-\operatorname{erf}\)
vanishes at infinity, computed cancellation-free so the tail keeps its
significance.}
\pair{special-functions/erf/6}{The imaginary error function
\(\operatorname{erfi}(z)=-i\,\operatorname{erf}(iz)\) has its own kernel.}
\pair{special-functions/erf/7}{The derivative is the Gaussian itself---so the whole
Taylor series is generated from this one rule.}
\pair{special-functions/erf/8}{The inverse function has exact endpoints:
\(\operatorname{erf}^{-1}(1)=\infty\).}
\pair{special-functions/erf/9}{And numeric values in between, on \([-1,1]\).}
\end{codepairs}

The split between \B{Erf} and \B{Erfc} (pairs~2 and~5) is not redundancy but
numerics: for large positive \(z\), \(\operatorname{erf}(z)\) is so close to \(1\)
that \(1-\operatorname{erf}(z)\) loses every significant digit to cancellation, so
\(\operatorname{erfc}\) is computed by its own cancellation-free
route.\calloutnote{Under the hood}{The kernels live in
\texttt{src/special\_functions/erf.c} (and \texttt{inverf.c} for the inverse). Real
arguments use the C library or MPFR's \texttt{erf}/\texttt{erfc}. A machine-precision
complex argument has its own kernel (\texttt{sf\_machine\_error.c}): Salzer's series,
with every square in \(e^{-x^2}\), \(e^{y^2}\) and the phase \(2xy\) formed exactly in
double-double arithmetic, since rounding \(x^2\) alone would cost nine ulps at \(x=3\)
before any algorithm ran. An arbitrary-precision complex argument sums the
cancellation-aware Maclaurin series of DLMF~7.6.2 in MPFR with \(|z|^2/\ln2\) guard
bits. Neither the C library nor MPFR ships an inverse error function, so
\B{InverseErf} is Newton's iteration on \(\operatorname{erf}(z)-s\), started from a
Winitzki seed and polished to the working precision.}

The same principle---name the antiderivative when it is not elementary---generates a
whole class of \emph{integral functions}, and Mathilda produces them directly when
you integrate.

\begin{codepairs}
\pair{special-functions/integrals/1}{The sine integral \(\operatorname{Si}\) is the
antiderivative of the sinc function, and \B{Integrate} returns it by name.}
\pair{special-functions/integrals/2}{The cosine integral \(\operatorname{Ci}\)
likewise.}
\pair{special-functions/integrals/3}{The exponential integral \(\operatorname{Ei}\)
names \(\int e^x/x\,dx\).}
\pair{special-functions/integrals/4}{And \(\int dx/\log x\) is the same function in
disguise---the logarithmic integral \(\operatorname{li}(x)=\operatorname{Ei}(\log x)\).}
\end{codepairs}

Each of these is a genuine function with the full three-faced treatment: it
evaluates, it differentiates, it has special values.

\begin{codepairs}
\pair{special-functions/integrals/5}{\B{SinIntegral} to twenty digits.}
\pair{special-functions/integrals/6}{\B{CosIntegral} to the same precision.}
\pair{special-functions/integrals/7}{\B{ExpIntegralEi} as well.}
\pair{special-functions/integrals/8}{And \B{LogIntegral}, whose growth rate is
the density of the primes.}
\pair{special-functions/integrals/9}{\(\operatorname{Si}\) saturates to \(\pi/2\) at
infinity---the Dirichlet integral \(\int_0^\infty \sin(x)/x\,dx\).}
\pair{special-functions/integrals/10}{The Fresnel integral \(C\) of optics and
diffraction is \B{FresnelC}.}
\pair{special-functions/integrals/11}{Its partner is \B{FresnelS}.}
\pair{special-functions/integrals/12}{And \(\operatorname{Si}\) differentiates back
to the cardinal sine \B{Sinc}---the loop closes.}
\end{codepairs}

That \(\int dx/\log x\) and \(\int e^x/x\,dx\) come back as the \emph{same} function
(pairs~3--4) is the logarithmic integral's claim to fame: \(\operatorname{li}(x)\)
is the leading term of the prime-counting function, and the substitution
\(t=\log x\) turns one integral into the other.\calloutnote{Under the hood}{At machine
precision the whole family runs on one set of double kernels
(\texttt{sf\_machine\_expint.c}), shared by the scalar, packed-array and \mcode{Compile}
paths: \B{ExpIntegralEi} on the real axis is Cody and Thacher's piecewise rational
fit (refitted against mpmath), and \mcode{SinIntegral}/\mcode{CosIntegral} follow Cephes. At
arbitrary precision \mcode{ExpIntegralEi} uses MPFR's correctly rounded
\texttt{mpfr\_eint} for \(x>0\) and an on-cut convergent series for \(x<0\); large
\(|z|\) switches to the asymptotic expansion, and the two regimes overlap and agree. \B{LogIntegral} is
literally \(\operatorname{Ei}(\log z)\), reusing the same kernel. The sine, cosine,
and Fresnel integrals sum their Maclaurin series for moderate arguments and their
trigonometric-prefactored asymptotic forms for large ones, with the complex plane
handled by the shared MPFR-complex toolkit.}

The exponential integral also comes in a graded family. The \emph{generalized
exponential integral} \(E_n(z)=\int_1^\infty e^{-zt}\,t^{-n}\,dt\), \B{ExpIntegralE},
is the Laplace transform of \(t^{-n}\) on \([1,\infty)\); integrating by parts lowers
\(n\) by one, so the whole ladder hangs from \(E_0(z)=e^{-z}/z\), and \(E_1\) is
\(\operatorname{Ei}\) reflected.\index{exponential integral!generalized} The order need
not be an integer.

\begin{codepairs}
\pair{special-functions/expinte/1}{The bottom rung is elementary.}
\pair{special-functions/expinte/2}{At the origin \(E_n(0)=1/(n-1)\) for \(n>1\).}
\pair{special-functions/expinte/3}{Thirty digits of \(E_2(1)\).}
\pair{special-functions/expinte/4}{A non-integer order, at machine precision.}
\pair{special-functions/expinte/5}{\(E_1(x)=-\operatorname{Ei}(-x)\) for \(x>0\): the
two heads cancel.}
\pair{special-functions/expinte/6}{Differentiation walks down the ladder,
\(E_n'=-E_{n-1}\).}
\end{codepairs}

\subsection{Bessel and Airy functions}
\label{ssec:sf-bessel}

Bessel functions are what solutions to the wave and heat equations look like in
cylindrical coordinates, and they are the archetype of a function defined by a
differential equation rather than a formula.\index{Bessel functions} \B{BesselJ}
is the function of the first kind, regular at the origin; \B{BesselY} the second
kind, singular there; and \B{BesselI}, \B{BesselK} the modified pair for the
hyperbolic version of the equation.

\begin{codepairs}
\pair{special-functions/bessel/1}{\(J_0(0)=1\); every higher \(J_n\) vanishes at the
origin.}
\pair{special-functions/bessel/2}{At half-integer order the Bessel function is
\emph{elementary}: \(J_{1/2}(x)=\sqrt{2/\pi x}\,\sin x\).}
\pair{special-functions/bessel/3}{And \(J_{-1/2}(x)=\sqrt{2/\pi x}\,\cos x\)---the
spherical Bessel functions in disguise.}
\pair{special-functions/bessel/4}{Numerics to any precision.}
\pair{special-functions/bessel/5}{The Frobenius series about the origin, generated on
demand.}
\pair{special-functions/bessel/6}{The recurrence shows itself in the derivative:
\(J_0'=-J_1\).}
\pair{special-functions/bessel/7}{The second kind \(Y_0\) is singular at the
origin.}
\pair{special-functions/bessel/8}{The modified \(I_0\) is evaluated at \(1\).}
\pair{special-functions/bessel/9}{And \(K_0\) is the exponentially decaying
modified solution.}
\end{codepairs}

The half-integer reductions (pairs~2--3) are a real theorem, not a lookup: they are
why a quantum-mechanics problem in a spherical well produces \(\sin\) and \(\cos\)
over \(\sqrt r\).\calloutnote{Theory}{The Bessel functions \(J_\nu\) and \(Y_\nu\)
are the two independent solutions of Bessel's equation
\(x^2y''+xy'+(x^2-\nu^2)y=0\). When \(\nu\) is a half-integer the equation degenerates
just enough for the solution to close in elementary terms: \(J_{1/2}\) and
\(J_{-1/2}\) are \(\sin\) and \(\cos\) scaled by \(\sqrt{2/\pi x}\), and the
three-term recurrence \(J_{\nu-1}(x)+J_{\nu+1}(x)=\tfrac{2\nu}{x}J_\nu(x)\) generates
every other half-integer order from them. Mathilda applies exactly this reduction
(\texttt{src/special\_functions/bessel.c}).}

\usagebox{BesselJ}

The Airy functions are the simplest functions that oscillate on one side of the
origin and decay on the other, the model for a wave near a turning
point.\index{Airy functions} \B{AiryAi} and \B{AiryBi} solve \(y''=xy\), and their
values at the origin tie the whole family back to the gamma function.

\begin{codepairs}
\pair{special-functions/airy/1}{At the origin \(\operatorname{Ai}\) closes
through the gamma function: \(\operatorname{Ai}(0)\) is \(1/(3^{2/3}\,\Gamma(2/3))\).}
\pair{special-functions/airy/2}{\(\operatorname{Bi}(0)\) shares the same
\(\Gamma(2/3)\), scaled differently.}
\pair{special-functions/airy/3}{The derivative at the origin brings in
\(\Gamma(1/3)\) instead.}
\pair{special-functions/airy/4}{\B{AiryAi} at the origin evaluates numerically, to
thirty digits.}
\pair{special-functions/airy/5}{At \(x=1\) it is already deep in the decaying
region.}
\pair{special-functions/airy/6}{\B{AiryAi} differentiates to its own derivative head,
\B{AiryAiPrime}, closing the family under \B{D}.}
\end{codepairs}

That \(\operatorname{Ai}(0)\) and \(\operatorname{Ai}'(0)\) involve \(\Gamma(2/3)\)
and \(\Gamma(1/3)\) respectively (pairs~1--3) is the thread tying this subsection
back to the first: the special functions are a single web, and the gamma function
sits near its center.\calloutnote{Theory}{The Airy equation \(y''=xy\) has no
singular points at finite \(x\), so both solutions are entire. The initial values
\(\operatorname{Ai}(0)=3^{-2/3}/\Gamma(2/3)\) and
\(\operatorname{Ai}'(0)=-3^{-1/3}/\Gamma(1/3)\) come from the Mellin-transform
representation of the solution, and every Taylor coefficient follows from the
equation by \(y^{(n+2)}=x\,y^{(n)}+n\,y^{(n-1)}\). The kernels are in
\texttt{src/special\_functions/airyai.c} and \texttt{airybi.c}.}

Two more members complete the Bessel neighborhood. Drive Bessel's equation with a
power, \(x^2y''+xy'+(x^2-\nu^2)y=4(x/2)^{\nu+1}/(\sqrt\pi\,\Gamma(\nu+\tfrac12))\), and
the particular solution is the \emph{Struve function} \B{StruveH},
\(\mathbf H_\nu(x)\), which turns up in the radiation impedance of a piston and in
unsteady aerodynamics.\index{Struve function} And the zeros of \(J_\nu\)---the
vibration frequencies of a drum, the cut-offs of a waveguide---have their own head,
\B{BesselJZero}, where \texttt{BesselJZero[nu, k]} is the \(k\)-th positive zero
\(j_{\nu,k}\).\index{Bessel functions!zeros}

\begin{codepairs}
\pair{special-functions/struve/1}{\(\mathbf H_\nu(0)=0\) for \(\operatorname{Re}\nu>-1\).}
\pair{special-functions/struve/2}{A machine value,}
\pair{special-functions/struve/3}{and thirty digits of \(\mathbf H_0(1)\).}
\pair{special-functions/struve/4}{For large \(x\), \(\mathbf H_0-Y_0\) decays like
\(2/(\pi x)\): the inhomogeneous term's footprint once the oscillations of \(Y_0\) are
subtracted---}
\pair{special-functions/struve/5}{---already within \(1\%\) at \(x=10\).}
\pair{special-functions/struve/6}{\(j_{0,1}\) has no closed form and stays a symbol,}
\pair{special-functions/struve/7}{but it is a number to any precision,}
\pair{special-functions/struve/8}{and \(J_0\) vanishes there to all thirty digits.}
\end{codepairs}

FLINT, which backs most of this section's arbitrary-precision numerics, ships no Struve
function, so Mathilda composes one.\calloutnote{Under the hood}{At arbitrary precision
\(\mathbf H_\nu(z)=(z/2)^{\nu+1}\,{}_1\tilde F_2(1;\tfrac32,\nu+\tfrac32;-z^2/4)\) is
evaluated in Arb, with the regularized \({}_1F_2\) from \texttt{acb\_hypgeom\_pfq} and
the working precision raised up front by about \(1.44|z|\) bits to pay for the
cancellation in an alternating series of size \(e^{|z|}\). At machine precision a
double kernel (\texttt{sf\_machine\_struve\_h}) picks between the convergent series and
the large-\(x\) form---Hankel's expansion of \(Y_\nu\) plus the asymptotic series
for \(\mathbf H_\nu-Y_\nu\) (DLMF~11.6.1), each cut at its smallest term---whichever
predicts the smaller error at that point; it serves the scalar call, packed arrays and \mcode{Compile}
alike. \mcode{BesselJZero} likewise gained a numeric kernel shared by all three
surfaces in v0.477.}

\subsection{Legendre functions of the first and second kind}
\label{ssec:sf-legendre}

Legendre's equation \((1-x^2)y''-2xy'+n(n+1)y=0\) is the angular part of Laplace's
equation on the sphere, and it too has two independent
solutions.\index{Legendre functions} \B{LegendreP} is the first: at a non-negative
integer degree it is the Legendre \emph{polynomial}, the workhorse of Gaussian
quadrature and multipole expansions. \B{LegendreQ} is the second kind---the other
solution of the same equation.\index{Legendre
functions!second kind}

\begin{codepairs}
\pair{special-functions/legendre/1}{The Legendre polynomials \(P_0\) through \(P_3\),
generated on demand.}
\pair{special-functions/legendre/2}{\(P_4\), with its characteristic rational
coefficients.}
\pair{special-functions/legendre/3}{Orthogonality on \([-1,1]\): the squared norm is
\(\int_{-1}^1 P_n^2 = 2/(2n{+}1)\), here \(2/5\).}
\pair{special-functions/legendre/4}{And distinct degrees are orthogonal---the
integral of \(P_2P_3\) is exactly zero.}
\pair{special-functions/legendre/5}{The second-kind solution \(Q_0\) is the
logarithm \(\tfrac12\log\frac{1+x}{1-x}=\operatorname{arctanh}x\).}
\pair{special-functions/legendre/6}{\(Q_1\) carries that same logarithm, times
\(x\), minus one.}
\pair{special-functions/legendre/7}{\(Q_2\) pairs the logarithm with the polynomial
\(P_2\)---the general pattern.}
\pair{special-functions/legendre/8}{And \(Q\) evaluates numerically like any other
function.}
\end{codepairs}

Pairs~3--4 are the property that makes the Legendre polynomials indispensable, and
the system computes them exactly rather than numerically.\calloutnote{Theory}{The
polynomials \(P_n\) are orthogonal on \([-1,1]\):
\(\int_{-1}^1 P_m(x)P_n(x)\,dx = \frac{2}{2n+1}\delta_{mn}\). This single relation is
what lets any function on the interval be expanded in a Legendre series, and it is the
foundation of Gauss--Legendre quadrature. The second-kind solution \(Q_n\) satisfies
the same differential equation but is singular at \(x=\pm1\)---the endpoints where the
logarithm in pairs~5--7 blows up---so a physical solution regular on the whole sphere
keeps only the \(P_n\) part. \B{LegendreQ} follows the pattern
\(Q_n(x)=P_n(x)\,Q_0(x)-W_{n-1}(x)\) with \(Q_0=\operatorname{arctanh}x\); the
implementation is in \texttt{src/special\_functions/legendre.c}.} The visible
structure of \(Q_2\) in pair~7---a Legendre polynomial multiplying
\(\operatorname{arctanh}x\), minus a lower-degree polynomial---is exactly that
recurrence made concrete.

\subsection{The classical orthogonal polynomials}
\label{ssec:sf-orthopoly}

The Legendre polynomials are the first of a small family. Change the interval or the
weight against which the polynomials are orthogonal and you get the other classical
families:\index{orthogonal polynomials} the Chebyshev polynomials \B{ChebyshevT} and
\B{ChebyshevU} on \([-1,1]\) with weights \((1-x^2)^{\mp1/2}\), the Hermite
polynomials \B{HermiteH} on the whole line with weight \(e^{-x^2}\) (the physicists'
normalization), and the Laguerre polynomials \B{LaguerreL} on \([0,\infty)\) with
weight \(x^a e^{-x}\). They are the eigenfunctions of the harmonic oscillator and the
hydrogen atom, the nodes of Gaussian quadrature, and---in the Chebyshev case---the
basis in which numerical analysts approximate everything else.

\begin{codepairs}
\pair{special-functions/orthopoly/1}{\(T_5\): odd degree, odd polynomial, leading
coefficient \(2^{4}\).}
\pair{special-functions/orthopoly/2}{\(U_4\), the second kind.}
\pair{special-functions/orthopoly/3}{\(H_5\), with leading coefficient \(2^5\).}
\pair{special-functions/orthopoly/4}{\(L_4\) has rational coefficients, alternating in
sign.}
\pair{special-functions/orthopoly/5}{The generalized \(L_n^{(a)}\) is a polynomial in
both \(x\) and the parameter.}
\pair{special-functions/orthopoly/6}{\(T_n(\cos t)=\cos nt\) is the definition the
others are measured against.}
\pair{special-functions/orthopoly/7}{Orthogonality under the Chebyshev weight,}
\pair{special-functions/orthopoly/8}{and the Hermite norm \(2^n n!\sqrt\pi\), here
\(48\sqrt\pi\).}
\end{codepairs}

What the families share is more useful than what separates
them.\calloutnote{Theory}{Any sequence of polynomials orthogonal against a positive
weight satisfies a \emph{three-term recurrence}\index{orthogonal
polynomials!three-term recurrence} \(p_{n+1}=(\alpha_nx+\beta_n)p_n-\gamma_np_{n-1}\)
(Favard's theorem is the converse). For the classical families the coefficients are
simple: \(T_{n+1}=2xT_n-T_{n-1}\), \(H_{n+1}=2xH_n-2nH_{n-1}\),
\((n+1)L_{n+1}=(2n+1-x)L_n-nL_{n-1}\). The same structure gives each family a
second-order differential equation and a derivative rule that stays inside the family,
which is what the derivative rules further on exploit.} The exact integer-order polynomials come from
FLINT (\texttt{fmpz\_poly\_chebyshev\_t}, \texttt{\_u}, \texttt{fmpz\_poly\_hermite\_h}),
with an integer recurrence as the fallback when Mathilda is built without it, and
Laguerre's rational coefficients from a rational recurrence.

\usagebox{ChebyshevT}

The recurrence matters even more for numbers than for symbols, and here the
orthogonal polynomials teach a lesson that applies to every special function: \emph{an
explicit formula is not an algorithm}. The expanded \(H_{100}\) has coefficients up to
about \(3\times10^{98}\), of alternating sign; evaluate it at \(x=5\) and terms as
large as \(2\times10^{118}\) cancel to leave a value of about
\(8\times10^{98}\)---twenty orders of magnitude below its largest term, so a double,
with its sixteen digits, has nothing left to say.\index{cancellation!polynomial evaluation}

\begin{codepairs}
\pair{special-functions/orthopoly-numeric/1}{A machine argument runs the three-term
recurrence in double precision.}
\pair{special-functions/orthopoly-numeric/2}{Substitution into the expanded polynomial
is wrong by a factor of about 6500---not a rounding error but a different number.}
\pair{special-functions/orthopoly-numeric/3}{The exact polynomial at the exact
argument, to twenty digits, sides with the recurrence.}
\pair{special-functions/orthopoly-numeric/4}{The same story for \(L_{100}(5)\),}
\pair{special-functions/orthopoly-numeric/5}{which substitution triples.}
\pair{special-functions/orthopoly-numeric/6}{Non-integer order: \(T_\nu(x)=\cos(\nu
\arccos x)\), here \(\cos(\tfrac12\arccos 0.3)\).}
\pair{special-functions/orthopoly-numeric/7}{Any precision, through the
hypergeometric form \({}_2F_1(-\nu,\nu;\tfrac12;\tfrac{1-x}{2})\).}
\pair{special-functions/orthopoly-numeric/8}{A complex order is no harder.}
\pair{special-functions/orthopoly-numeric/9}{\(T_n'=nU_{n-1}\) links the two Chebyshev
kinds,}
\pair{special-functions/orthopoly-numeric/10}{\(H_n'=2nH_{n-1}\) drops the
degree by one,}
\pair{special-functions/orthopoly-numeric/11}{and \(L_n'=-L_{n-1}^{(1)}\) raises the
Laguerre parameter.}
\pair{special-functions/orthopoly-numeric/12}{The associated Legendre function
\(P_3^1\) carries the factor \(\sqrt{1-x^2}\).}
\pair{special-functions/orthopoly-numeric/13}{\(P_{30}^4\) over a list of machine
reals.}
\end{codepairs}

\calloutnote{Pitfall}{\texttt{HermiteH[100, x] /. x -> 5.} (pair~2) is the natural way
to ``plug in a number'', and it is the wrong one. \texttt{/.} substitutes into the
expression \emph{after} it has been expanded, so the evaluation order is fixed by the
monomial basis, the worst-conditioned one available. Pass the number to the function
itself, \texttt{HermiteH[100, 5.]}, and Mathilda chooses the stable algorithm. The same
holds for every polynomial family in this section, \mcode{LegendreP} included.} The
non-integer orders (pairs~6--8) leave the polynomial world altogether: there each
family is a hypergeometric function---\(T_\nu\) and \(U_\nu\) a \({}_2F_1\), \(H_\nu\)
and \(L_\nu^{(a)}\) a confluent \({}_1F_1\)---and is evaluated as one, at machine or
arbitrary precision, real or complex.\calloutnote{Performance}{Each of the four families, and
\mcode{LegendreP[n, x]} itself, has one double-precision recurrence kernel (with a
complex arm) shared by the scalar call, the packed-array and \B{NDArray} maps and
\mcode{Compile}, so those surfaces agree bit for bit (\S\ref{ssec:sf-machine}). The
associated Legendre functions over a machine array are evaluated by Horner's rule in
\texttt{double} rather than as a sum of monomials, each of which used to cost a
full-buffer \texttt{pow}; that change (v0.376)
took \texttt{LegendreP[30, 4, xs]} over a million points from not finishing in
85~seconds to about a tenth of a second on this build.}

\subsection{Legendre elliptic integrals}
\label{ssec:sf-elliptic}

Measure the arc of an ellipse, or the period of a pendulum that swings more than a
few degrees, and the integral you get cannot be done in elementary
terms.\index{elliptic integrals} That is not a failure of technique: Liouville's
theory says no elementary antiderivative exists. Legendre's response was to name
three canonical integrals and reduce every other one to
them,\index{elliptic integrals!Legendre normal form} and those three are
\B{EllipticK} and \B{EllipticF} (first kind, complete and incomplete), \B{EllipticE}
(second kind), and \B{EllipticPi} (third kind):
\[
  F(\phi\,|\,m)=\int_0^\phi\frac{dt}{\sqrt{1-m\sin^2 t}},\qquad
  E(\phi\,|\,m)=\int_0^\phi\!\sqrt{1-m\sin^2 t}\;dt,
\]
\[
  \Pi(n;\phi\,|\,m)=\int_0^\phi\frac{dt}{(1-n\sin^2 t)\sqrt{1-m\sin^2 t}},
\]
with the \emph{complete} integrals \(K(m)=F(\pi/2\,|\,m)\), \(E(m)=E(\pi/2\,|\,m)\)
and \(\Pi(n\,|\,m)\) taking the upper limit to \(\pi/2\). One warning before any
example: the argument \(m\) is the \textbf{parameter} \(m=k^2\), not the modulus
\(k\)---Mathilda follows the Wolfram convention, and this is the one place a reader
of these signatures goes silently
wrong.\index{elliptic integrals!parameter versus modulus}

\begin{codepairs}
\pair{special-functions/elliptic/1}{With \(m=0\) the integrand is \(1\), so \(K\) is
the quarter period \(\pi/2\).}
\pair{special-functions/elliptic/2}{At \(m=1/2\) a closed form exists---one of the
lemniscatic singular values, through \(\Gamma\).}
\pair{special-functions/elliptic/3}{The same number to thirty digits, through Arb's
ball arithmetic.}
\pair{special-functions/elliptic/4}{At \(m=1\) the integrand loses its last
denominator at the upper limit and \(K\) diverges.}
\pair{special-functions/elliptic/5}{The expansion at \(m=0\), from a dedicated
\B{Series} kernel.}
\pair{special-functions/elliptic/6}{A full quarter period closes the incomplete form
into the complete one.}
\pair{special-functions/elliptic/7}{The amplitude derivative \emph{is} the
integrand---the fundamental theorem of calculus, made visible.}
\pair{special-functions/elliptic/8}{The quasi-period
\(F(\phi+k\pi\,|\,m)=F(\phi\,|\,m)+2kK(m)\), confirmed as a residual.}
\pair{special-functions/elliptic/9}{And the parameter derivative of the second kind,
expressed in \(E\) and \(K\) themselves.}
\end{codepairs}

Pair~5 is a case where a general method genuinely cannot work and a special one must.
The generic Taylor route would differentiate \(K\) and evaluate at \(m=0\), but
\(dK/dm=(E-(1-m)K)/(2m(1-m))\) is \((\pi/2-\pi/2)/0\) there---indeterminate---so the
expansion was abandoned before it produced even its leading term. The closed
coefficient recurrence \(a_k=a_{k-1}((2k-1)/2k)^2\) has no such trouble.

\usagebox{EllipticK}

The next four examples are the family's two genuine traps, and they are worth more
than their line count. The first: \(E(\phi\,|\,1)=\int_0^\phi|\cos t|\,dt\), which
equals \(\sin\phi\) only while \(|\phi|\le\pi/2\). Applied without that guard it
reports \(\sin 2=0.909\) for a true \(2-\sin 2=1.091\)---a wrong answer that no
amount of precision would reveal, since the formula is exactly right on the strip
where most examples live. The second: for \(n>1\) the path of \(\Pi\) crosses the
pole at \(\sin^2 t=1/n\), and the value beyond it is genuinely
complex\index{elliptic integrals!third kind}---not the real Cauchy principal value it
is often said to be.

\begin{codepairs}
\pair{special-functions/elliptic/10}{On the principal strip the \(m=1\) identity
holds, and \(E(1/2\,|\,1)=\sin(1/2)\) exactly.}
\pair{special-functions/elliptic/11}{Beyond it the identity is false; the engine
declines to apply it and \(N\) recovers the true \(2-\sin 2\).}
\pair{special-functions/elliptic/12}{A vanishing characteristic turns the third kind
into the first.}
\pair{special-functions/elliptic/13}{Past the pole, \(\Pi(3/2\,|\,1/2)\) is complex,
and both digits and sign matter.}
\pair{special-functions/elliptic/14}{The perimeter of an ellipse with \(a=1\) and
\(e^2=1/4\)---the problem the second kind was named for.}
\pair{special-functions/elliptic/15}{A pendulum released at \(30^\circ\) swings
\(1.74\%\) slower than the small-angle law predicts.}
\pair{special-functions/elliptic/16}{And \(K\) is a Gauss hypergeometric function in
disguise: \(K(m)=\tfrac{\pi}{2}\,{}_2F_1(\tfrac12,\tfrac12;1;m)\).}
\end{codepairs}

Pairs~14--15 are the two problems that put these integrals in every textbook, and
pair~16 is the bridge to the next subsection.\calloutnote{Under the hood}{Numerically
the integrals are computed by \emph{Carlson's symmetric forms}---\(R_F\), \(R_D\),
\(R_J\) and \(R_C\), evaluated by duplication until the arguments agree to a
tolerance and then closed with a fifth-order
tail.\index{Carlson symmetric forms} The arrangement rewards care: \(R_F\) and
\(R_D\) walk an identical argument sequence, so computing them in one loop buys three
square roots per step instead of six, and the first argument of \(R_F\) must be
spelled \(\cos^2 r\) rather than \(1-\sin^2 r\)---near \(r=\pi/2\) the subtraction's
absolute error is the size of the answer, which once cost eight digits of
\(F(\pi/2-10^{-8}\,|\,0.99)\). The machine kernels cover the real principal domain at
a few ulp and \emph{decline} outside it, where FLINT's \texttt{acb\_elliptic\_*}
answers in rigorous ball arithmetic on the whole complex plane. Everything is in
\texttt{src/special\_functions/elliptic.c}.}

\subsection{The hypergeometric family: one function to rule them all}
\label{ssec:sf-hypergeometric}

Almost every function in this section is a special case of one
master series.\index{hypergeometric function} The generalized hypergeometric
function
\[
  {}_pF_q\!\left(\begin{matrix}a_1,\dots,a_p\\ b_1,\dots,b_q\end{matrix};z\right)
  = \sum_{k=0}^\infty
    \frac{(a_1)_k\cdots(a_p)_k}{(b_1)_k\cdots(b_q)_k}\,\frac{z^k}{k!}
\]
is a power series whose term ratio is a fixed rational function of the summation
index---and that one structural constraint is general enough to contain the
exponential, the logarithm, the trigonometric functions, the Bessel functions, the
error function, the Legendre functions, and the elliptic integrals of
\S\ref{ssec:sf-elliptic} as special cases. Mathilda exposes it as
\B{HypergeometricPFQ}, with the classical low-order cases \B{Hypergeometric2F1}
(Gauss), \B{Hypergeometric1F1} (Kummer's confluent), and \B{Hypergeometric0F1}
(the confluent limit) as named heads. The system's job is to \emph{recognize} when a
hypergeometric collapses to something elementary, and it does.

\begin{codepairs}
\pair{special-functions/hypergeometric/1}{Gauss's \({}_2F_1(1,1;2;x)\) is a
logarithm in disguise: \(-\log(1-x)/x\).}
\pair{special-functions/hypergeometric/2}{A negative-integer parameter terminates the
series into a polynomial---here \((1-x)^3\).}
\pair{special-functions/hypergeometric/3}{Kummer's \({}_1F_1(1;2;x)\) is
\((e^x-1)/x\).}
\pair{special-functions/hypergeometric/4}{When its two parameters agree it
is the exponential itself.}
\pair{special-functions/hypergeometric/5}{The confluent \({}_0F_1(;1/2;x)\) is a
hyperbolic cosine of a square root.}
\pair{special-functions/hypergeometric/6}{The generalized \({}_pF_q\) carries any
number of parameters, and reduces by the same rules---here matching pair~1.}
\pair{special-functions/hypergeometric/7}{Where no elementary form exists, it is a
number: this one is \(2\log 2\).}
\end{codepairs}

The unification is not just aesthetic; it is how a CAS avoids re-implementing the
same reductions a dozen times.\calloutnote{Theory}{A function is
\emph{hypergeometric} when the ratio of consecutive series coefficients,
\(c_{k+1}/c_k\), is a rational function of \(k\)---and astonishingly many functions
are. The exponential is \({}_0F_0\); \(\log(1-x)=-x\,{}_2F_1(1,1;2;x)\) (pair~1);
\(J_\nu(x)\) is a \({}_0F_1\) up to a power factor (compare pair~5's \(\cosh\), which
is \(J_{-1/2}\) territory); \(\operatorname{erf}\), the Legendre and Gegenbauer
polynomials, and the Airy functions are all \({}_pF_q\) with particular parameters.
The two levers that produce a closed form are visible above: a \emph{negative
integer} upper parameter truncates the infinite series to a polynomial (pair~2), and
special parameter relations trigger classical summation theorems (pairs~1, 3--5).
Mathilda's reductions live in \texttt{src/special\_functions/hypergeopfq.c}.}

Three functions extend the reach of the elementary world one step further and round
out the family. The polylogarithm \B{PolyLog} generalizes the logarithm to a whole
ladder of orders;\index{polylogarithm} the Lerch transcendent \B{LerchPhi} contains
the polylogarithm, the Hurwitz zeta, and much else as special
cases;\index{Lerch transcendent} and the Lambert \(W\) function \B{ProductLog}
inverts \(w\,e^w=z\), the transcendental equation that appears whenever an unknown
sits both inside and outside an exponential.\index{Lambert W function}

\begin{codepairs}
\pair{special-functions/polylog/1}{Order one is the logarithm:
\(\operatorname{Li}_1(x)=-\log(1-x)\).}
\pair{special-functions/polylog/2}{The dilogarithm at \(1\) is \(\zeta(2)=\pi^2/6\).}
\pair{special-functions/polylog/3}{At \(-1\) it is \(-\pi^2/12\).}
\pair{special-functions/polylog/4}{And the celebrated closed form
\(\operatorname{Li}_2(1/2)=\pi^2/12-\tfrac12\log^2 2\).}
\pair{special-functions/polylog/5}{At argument \(1\) the order-\(s\) polylog is
\(\zeta(s)\), so \(\operatorname{Li}_3(1)=\zeta(3)\).}
\pair{special-functions/polylog/6}{The Lerch transcendent reduces to the same
logarithm for these arguments---the master function containing all the rest.}
\pair{special-functions/polylog/7}{The Lambert \(W\): \(W(e)=1\), since
\(1\cdot e^1=e\).}
\pair{special-functions/polylog/8}{And \(W(-1/e)=-1\), the branch point.}
\pair{special-functions/polylog/9}{Its value at \(1\) is the omega constant,
\(0.5671\ldots\).}
\pair{special-functions/polylog/10}{Its derivative is written implicitly in terms of
itself---the signature of a function defined by inversion.}
\end{codepairs}

The dilogarithm identities (pairs~2--4) are among the most beautiful closed forms in
analysis, and that Mathilda knows \(\operatorname{Li}_2(1/2)=\pi^2/12-\tfrac12\log^2
2\) exactly rather than as a decimal is the whole point of a symbolic
system.\calloutnote{Under the hood}{\B{PolyLog}
(\texttt{src/special\_functions/polylog.c}), \B{LerchPhi} (\texttt{lerchphi.c}), and
\B{ProductLog} (\texttt{productlog.c}) each carry their own special-value tables and
their own numeric kernels. The Lerch transcendent
\(\Phi(z,s,a)=\sum_{k\ge0}z^k/(k+a)^s\) is the genuine master function of the three:
\(\operatorname{Li}_s(z)=z\,\Phi(z,s,1)\) and \(\zeta(s,a)=\Phi(1,s,a)\), so pair~6's
reduction and the Hurwitz zeta of \S\ref{ssec:sf-zeta} are two faces of one series.
\B{ProductLog}'s implicit derivative (pair~10) comes from differentiating
\(w\,e^w=z\) directly, since no explicit formula for \(W\) exists.}

The polylogarithm has a two-index generalization that appears wherever Feynman
integrals or iterated integrals of logarithms are done by hand. \emph{Nielsen's
generalized polylogarithm}\index{polylogarithm!Nielsen generalized}
\[
  S_{n,p}(z)=\frac{(-1)^{n+p-1}}{(n-1)!\,p!}\int_0^1\frac{\log^{n-1}t\,\log^p(1-zt)}{t}\,dt
\]
is \texttt{PolyLog[n, p, z]}, with \(S_{n-1,1}=\operatorname{Li}_n\). Its exact special
values are not built in---\texttt{PolyLog[1, 2, 1/2]} stays as it is---but its numerics
are, at any precision.

\begin{codepairs}
\pair{special-functions/nielsen/1}{\(S_{1,2}(\tfrac12)\) at machine precision,}
\pair{special-functions/nielsen/2}{to thirty digits,}
\pair{special-functions/nielsen/3}{and the classical closed form
\(\zeta(3)/8-\tfrac16\log^32\) it equals.}
\pair{special-functions/nielsen/4}{\(S_{2,2}(1)\) to thirty digits.}
\pair{special-functions/nielsen/5}{\(\zeta(4)/4\) gives the same digits, so
\(S_{2,2}(1)=\zeta(4)/4\).}
\pair{special-functions/nielsen/6}{Past \(z=1\), on the cut, integer \(n\) and \(p\)
continue analytically.}
\pair{special-functions/nielsen/7}{A non-integer order off the unit disk has no agreed
continuation, and stays unevaluated.}
\pair{special-functions/nielsen/8}{Over a list of machine reals.}
\pair{special-functions/nielsen/9}{\mcode{LerchPhi} with \mcode{DoublyInfinite -> True}
sums over all \(k\in\mathbb Z\): \(\sum_k(k+\tfrac12)^{-2}=\pi^2\).}
\end{codepairs}

Pair~7 is deliberate: Mathilda evaluates \(S_{n,p}\) on exactly the domain where the
reference system does---any \(n\) and any \(p\) with \(\operatorname{Re}p>0\) on the closed unit disk (\(z\ne1\)), \(\operatorname{Re}n>0\)
at \(z=1\), integer \(n\le0\) anywhere through the closed form \((z\,d/dz)^{-n}\), and
integer \(n,p\) beyond the disk---and leaves the rest alone rather than choosing a branch
nobody else would.\calloutnote{Under the hood}{In Arb at any precision \(S_{n,p}\) is
summed by a Stirling-number series for integer \(p\) and small \(|z|\), by a
generalized power series built with Miller's recurrence for any \(p\), and otherwise
by tanh--sinh quadrature of the integral above, split at \(t=1/z\) when the path meets
the cut. A machine kernel (v0.479) runs the same methods in double precision and
declines to Arb when its error estimate passes \(2\times10^{-14}\); over a 479-value
test grid it is within \(6.1\times10^{-15}\) of 30-digit values.}

\subsection{Kummer, Tricomi and Whittaker: the confluent family}
\label{ssec:sf-whittaker}

Kummer's equation \(zw''+(b-z)w'-aw=0\) is the confluent limit of Gauss's
hypergeometric equation, and like every second-order equation it has two independent
solutions.\index{confluent hypergeometric functions} The one regular at the origin is
\({}_1F_1(a;b;z)\), already met as \mcode{Hypergeometric1F1}; the other, singled out
by its algebraic decay \(U\sim z^{-a}\) as \(z\to+\infty\), is Tricomi's function
\B{HypergeometricU}. Remove the first-derivative term and Kummer's equation becomes
Whittaker's, \(w''+\bigl(-\tfrac14+\tfrac kz+\tfrac{1/4-m^2}{z^2}\bigr)w=0\), whose
solutions \B{WhittakerM} and \B{WhittakerW} are the same two functions dressed in a
prefactor,
\[
  M_{k,m}(z)=e^{-z/2}z^{m+1/2}\,{}_1F_1\bigl(m-k+\tfrac12;2m+1;z\bigr),\qquad
  W_{k,m}(z)=e^{-z/2}z^{m+1/2}\,U\bigl(m-k+\tfrac12,2m+1,z\bigr).
\]
The radial wave functions of the hydrogen atom, the Coulomb wave functions, the
parabolic cylinder functions and the incomplete gamma function are all Whittaker
functions in disguise.\index{Whittaker functions}

\begin{codepairs}
\pair{special-functions/whittaker/1}{When \(b=a+1\), \(U\) is exactly its own
asymptotic form \(z^{-a}\).}
\pair{special-functions/whittaker/2}{So \(U(2,3,z)=z^{-2}\).}
\pair{special-functions/whittaker/3}{Differentiation shifts both parameters,
\(U'=-a\,U(a{+}1,b{+}1,z)\).}
\pair{special-functions/whittaker/4}{Thirty digits of \(U(\tfrac13,2,5)\).}
\pair{special-functions/whittaker/5}{\(M_{k,m}\) at machine precision,}
\pair{special-functions/whittaker/6}{\(W_{k,m}\) at the same point,}
\pair{special-functions/whittaker/7}{and \(W\) with a complex index.}
\pair{special-functions/whittaker/8}{At \(m=\tfrac12\) the origin value of \(W\) is
finite: \(1/\Gamma(1-k)\).}
\pair{special-functions/whittaker/9}{The derivative closes in \(W\) itself, with the
index \(k\) raised by one.}
\pair{special-functions/whittaker/10}{The defining identity of \(M\), checked at a
machine point.}
\end{codepairs}

\usagebox{WhittakerW}

Tricomi's \(U\) is the harder of the two to compute, because no single expansion
serves the whole plane.\calloutnote{Under the hood}{The machine kernel
(\texttt{sf\_machine\_hyperu}) switches between three representations: the asymptotic
\({}_2F_0\) series for large \(z\); for integer \(b\), the logarithmic series of
DLMF~13.2.9, since there the two \({}_1F_1\) solutions coincide and their difference
must be taken analytically; and for non-integer \(b\), Kummer's connection formula
\(U=\frac{\Gamma(1-b)}{\Gamma(a-b+1)}M(a,b,z)+\frac{\Gamma(b-1)}{\Gamma(a)}z^{1-b}M(a-b+1,2-b,z)\).
A guard measures the cancellation in that difference and declines to the rigorous path,
FLINT's \texttt{acb\_hypgeom\_u}, which also answers at arbitrary precision and on
\(z\le0\). \mcode{WhittakerW} multiplies the same kernel by its prefactor (folding
\(m\to|m|\), since \(W_{k,m}=W_{k,-m}\)), so it inherits the accuracy and the
array and \mcode{Compile} surfaces of \(U\) unchanged.}

\subsection{The Meijer \texorpdfstring{\(G\)}{G}-function}
\label{ssec:sf-meijerg}

The generalized hypergeometric \({}_pF_q\) is a power series; Meijer's \(G\) is the
contour integral that contains it, and it is the largest family in common use that is
still closed under the operations a CAS performs on integrands.\index{Meijer
G-function} With parameter lists split as \(a_1,\dots,a_n\mid a_{n+1},\dots,a_p\) and
\(b_1,\dots,b_m\mid b_{m+1},\dots,b_q\),
\[
  G^{m,n}_{p,q}\!\left(z\,\middle|\,\begin{matrix}a_1,\dots,a_p\\
  b_1,\dots,b_q\end{matrix}\right)
  =\frac1{2\pi i}\int_L
  \frac{\prod_{j=1}^m\Gamma(b_j-s)\prod_{j=1}^n\Gamma(1-a_j+s)}
       {\prod_{j=m+1}^q\Gamma(1-b_j+s)\prod_{j=n+1}^p\Gamma(a_j-s)}\,z^s\,ds,
\]
written \texttt{MeijerG[\{\{a1..an\}, \{a(n+1)..ap\}\}, \{\{b1..bm\}, \{b(m+1)..bq\}\}, z]}.
Closing the contour around the poles of \(\Gamma(b_j-s)\) and summing residues turns
the integral into \(m\) hypergeometric series---\emph{Slater's
theorem}\index{Slater's theorem}---and that is how \B{MeijerG} recognizes the named
functions it contains.

\begin{codepairs}
\pair{special-functions/meijerg/1}{The simplest \(G\) of all, \(G^{1,0}_{0,1}(z|0)\), is
\(e^{-z}\).}
\pair{special-functions/meijerg/2}{One \(a\) and one \(b\) give a binomial power.}
\pair{special-functions/meijerg/3}{\(G^{1,2}_{2,2}\) with these parameters is
\(\log(1+z)\).}
\pair{special-functions/meijerg/4}{\(G^{1,0}_{0,2}\) is a Bessel \(J\) of \(2\sqrt z\)---}
\pair{special-functions/meijerg/5}{---and \(G^{2,0}_{0,2}\) a Bessel \(K\).}
\pair{special-functions/meijerg/6}{Three \(b\)'s have no single classical name; the
parameters come back sorted into canonical order.}
\pair{special-functions/meijerg/7}{\B{FunctionExpand} writes out the Slater sum: three
\({}_0F_2\) terms with gamma-function weights.}
\pair{special-functions/meijerg/8}{Machine values, threaded over a list,}
\pair{special-functions/meijerg/9}{and the same point to twenty-five digits.}
\pair{special-functions/meijerg/10}{Here \(b_1,b_2\) differ by an integer---the log
case---so no Slater expansion exists and the symbolic form stays unevaluated.}
\pair{special-functions/meijerg/11}{The number still exists: thirty digits at
\(z=3\).}
\end{codepairs}

The automatic reductions (pairs~1--5) and \mcode{FunctionExpand} (pair~7) are two
strengths of the same rewrite: \mcode{MeijerG} applies Slater's theorem only when the
result is a classical function, while \mcode{FunctionExpand} applies it whenever an
expansion exists. That is, at present, \mcode{FunctionExpand}'s only rule; every other
expression passes through it unchanged.\calloutnote{Theory}{When two \(b_j\) among the
first \(m\) differ by an integer, the poles of the corresponding \(\Gamma(b_j-s)\)
collide, the residues become derivatives, and the expansion acquires powers of
\(\log z\)---the same mechanism that puts a logarithm into \(Y_0\) and \(K_0\). This
\emph{log case} has no finite sum of \({}_pF_q\)'s, which is why pairs~10--11 stay
symbolic and evaluate only as numbers. The reason \(G\) matters is closure: the
product of two \(G\)'s integrated over \((0,\infty)\) is again a \(G\), so an integrator
that rewrites both factors as \(G\)-functions can do whole classes of definite
integrals by parameter bookkeeping, the Mellin-transform route of
\S\ref{ssec:calc-definite}.}\calloutnote{Performance}{Numerically a parameter set is turned
into a \emph{plan} once---the \(z\)-independent \(\pi/\sin\) and gamma products of every
Slater term---and each point costs one hypergeometric kernel call per term; a
cancelling sum is re-evaluated in Arb. The log case has its own machine kernel
(v0.446), the residue series \(\sum_k w^{\beta+k}\sum_j c_{kj}\log^jw\) with the
coefficients tabulated once per plan and summed in double-double where it cancels: a
scalar call went from \(456\,\mu\)s to about \(3\,\mu\)s, a \(10^6\)-element buffer from
119~s to 0.3~s. In the special-functions benchmark (\S\ref{ssec:sf-machine}) Mathilda is
the fastest engine in all nine \mcode{MeijerG} cells.}

\usagebox{MeijerG}

\subsection{Machine precision: one kernel, five surfaces}
\label{ssec:sf-machine}

Two rules govern what happens when a special function meets a floating-point
number.\index{special functions!machine precision} The first is about mixed input: a
call with one inexact argument and other exact but irrational ones is numericalized as
a whole, at the lowest precision present, rather than left half-evaluated. The second
is about where the number is computed, and it is the subject of the rest of this
subsection.

\begin{codepairs}
\pair{special-functions/mixed/1}{A machine \(s\) numericalizes the exact
\(a=\sqrt2\).}
\pair{special-functions/mixed/2}{A 30-digit \(s\) gives a 30-digit answer.}
\pair{special-functions/mixed/3}{The same rule for an irrational order.}
\pair{special-functions/mixed/4}{The \(q\)-Pochhammer symbol \B{QPochhammer},
\((a;q)_n=\prod_{k<n}(1-aq^k)\), exactly,}
\pair{special-functions/mixed/5}{and at a non-integer \(n\), where it is a ratio of
infinite products.}
\end{codepairs}

A machine-precision special function can be reached five ways: as a scalar call, an
element of an unpacked list, a packed array, a visible \mcode{NDArray}, or inside
\mcode{Compile}d code (\S\ref{sec:ds-packed}, \S\ref{ch:compilation}). In Mathilda all
five call the \emph{same} raw-double kernel, so they return the same bits, not merely
the same digits.\index{special functions!bit-for-bit surfaces} Since v0.479 that holds
for the \(n\)-ary heads too---\mcode{LerchPhi}, \mcode{PolyLog}, the \({}_1F_1\) and
\({}_2F_1\), \mcode{HypergeometricU}, \mcode{WhittakerM}/\mcode{W},
\mcode{EllipticPi}, \mcode{LaguerreL}, \mcode{DirichletL}---with the array in any one
operand.

\begin{codepairs}
\pair{special-functions/surfaces/1}{Twenty machine values of the index \(k\).}
\pair{special-functions/surfaces/2}{The array sits in the \emph{first} operand of a
three-argument head,}
\pair{special-functions/surfaces/3}{and the result is still a packed array, by
\B{PackedArrayQ}.}
\pair{special-functions/surfaces/4}{A scalar call on one element gives identical
bits (\B{SameQ}, not \mcode{Equal}),}
\pair{special-functions/surfaces/5}{so does compiled code,}
\pair{special-functions/surfaces/6}{with the array in the same slot,}
\pair{special-functions/surfaces/7}{and so does a visible \mcode{NDArray}.}
\pair{special-functions/surfaces/8}{\B{CompileDiagnostics} confirms the head lowered:
no fall-back to the interpreter.}
\end{codepairs}

Bit-for-bit agreement is a correctness property, not a nicety: a numerical method that
evaluates the same function by two routes and differences the results (a Newton step,
a finite difference, a convergence test) sees the disagreement between the routes as
signal. It is also the reason the scalar path is fast.\calloutnote{Performance}{The
four-engine special-functions benchmark in
\texttt{benchmarks/97-special-functions/} (its report is
\texttt{special\_function\_performance.tex}) times 58 rows---every numerically
evaluable special function, some in several variants---against SciPy (with mpmath for
arbitrary precision), Maxima and the reference system, with distinct arguments per
call so no engine can serve a cached value. In the report's latest run (v0.446, with the
\mcode{DirichletL} rows from v0.450) Mathilda has the lowest
per-call latency in \emph{every} machine-real and machine-complex row, with medians of
\(8.9\times\) and \(10.2\times\) below the fastest other engine, and its
\(10^6\)-element packed buffers beat SciPy for every head SciPy implements. Every
machine kernel was gated against mpmath at a median of about one ulp before it was
wired in. At 50 digits the picture is mixed: a median of \(2.7\times\) ahead, but
fastest in 46 of 57 rows, with Maxima an order of magnitude quicker on \mcode{BesselY}
and \mcode{BesselK}. The machine regimes were closed (v0.436--v0.444) by routing every
scalar call through one gate to the kernel the array paths already used, and by
writing native double re/im complex kernels in place of MPFR pairs.}

\subsection*{Where this connects}

The special functions are the point where the three layers of Mathilda meet on a
single object. Their \emph{exact} values are computed by the same rational and
symbolic arithmetic as \S\ref{sec:arithmetic} and \S\ref{sec:algebra}; their
\emph{identities}---the derivatives that close the gamma, Bessel, Airy, Legendre and
elliptic families under \B{D}, the series expansions, the antiderivatives that
\B{Integrate} returns by name---are what let them take part in the calculus of
\S\ref{sec:calculus} rather than merely being evaluated at its end; and their
\emph{numerics}, backed by MPFR on the real line and by FLINT's rigorous Arb ball
arithmetic in the complex plane, are what the numerical methods of
\S\ref{sec:numcalc} call when an integral or a differential equation resolves into
one of these functions. The orthogonality of the Legendre polynomials underwrites
Gaussian quadrature; the gamma and zeta functions, and the Dirichlet
\(L\)-functions that twist them, thread through analytic number theory; and the
hypergeometric family, with Meijer's \(G\) at its summit, is the organizing idea that keeps the whole
collection from being a bag of unrelated formulas. For the exhaustive options and
corner cases of any function named here---the branch-cut conventions, the incomplete
and generalized forms, the full list of recognized reductions---its reference page
is one \texttt{?Name} away at the prompt.
